\input{./._relpath-to-latexroot.ltx}
\documentclass[\latexroot/\projectname]{subfiles}
\input{\latexroot/@resources/subfile-setup.ltx}

\begin{document}

\section{Introduction}
\whenintegrated{\label{sec:intro}} % integrated doc: SST for labels
\setcounter{page}{0}\pagenumbering{arabic}

If the government purchases $\$ 1$ worth of goods, by how much does private consumption increase or decrease? Although the question of the consumption response to government spending is very old, the literature still lacks consensus. For example, Ramey and Shapiro~\cite{Ramey1998-as} find that exogenous increases in defense spending decrease private consumption. On the other hand, Blanchard and Perotti (2002) and Gali et al.~\cite{Gali2007-ji} find that exogenous fiscal expansions increase private

The editor in charge of this paper was Dirk Krueger.\\
consumption.\footnote{In a review of the literature and empirical methods, Hall~\cite{Hall2009-mz} finds a consumption multiplier from somewhat negative to 0.5. Other recent contributions include Ramey (2011), Auerbach and Gorodnichenko (2012), Ramey~\cite{Ramey2016-tn}, and Ramey and Zubairy~\cite{Ramey2018-rc}.} This disconnect is worrisome since consumer spending is the largest component of national income and its response is a key determinant of the fiscal multiplier.

In this paper, we first estimate the response of consumer spending to fiscal stimulus. In particular, we use regional variation in the spending component of the American Recovery and Reinvestment Act (ARRA) to estimate the local effect of government spending on consumer spending. The local effect is the relative change in regional outcomes in response to a relative change in government spending across regions.\footnote{Nakamura and Steinsson~\cite{Nakamura2011-dw} refer to this object, in a fiscal policy context, as the open economy relative multiplier.}

Second, we translate the local fiscal multiplier to an aggregate fiscal multiplier using a multi-region, new Keynesian model with heterogeneous agents and incomplete markets. This is because estimates based on regional variation are not always representative of aggregate effects. Such estimates ignore general equilibrium effects that cannot be separately identified in cross-regional regressions (~\cite{Nakamura2018-wz}; ~\cite{Chodorow-Reich2019-zl}).

The ARRA, implemented between 2009 and 2012, was a very large programme by historical standards. The spending component of the Act allocated roughly $\$ 228$ billion. Our consumer spending data come from two separate sources. First, we collect store-level information on retail purchases from the Nielsen retail scanner data. Second, we construct individual-level spending on vehicles by measuring changes in auto credit balances from the FRB NY Consumer Credit Panel/Equifax.

We estimate that a $\$ 1$ increase in ARRA spending within a county increases local nondurable spending by $\$ 0.29$ and local auto spending by $\$ 0.09$. We address potential endogeneity of government spending in two ways. First, we show that ARRA funding did not significantly target low-income areas. Second, we use a narrative instrumental variable approach. In particular, we identify components of ARRA funding that did not explicitly target local economic recovery.

We translate our estimated local fiscal multiplier to an aggregate fiscal multiplier using a general equilibrium, two region, new Keynesian model. In the model, each region produces a final good, which is purchased by the local consumers as well as the government. The final good is produced using both local and foreign intermediate inputs. Due to home bias, final goods are produced using a larger share of local inputs. Trade linkages-expressed in the degree of home bias-provide one channel for government spending to affect economic activity across regions. The government maintains a fiscal union because it finances spending using federal taxes. The monetary authority targets a common nominal interest rate for all regions (currency union).

Our multi-region model is novel in an important dimension: each region is populated by heterogeneous households who face incomplete asset markets (~\cite{Huggett1993-ft}; ~\cite{Aiyagari1994-qp}). Heterogeneity in our model takes the form of idiosyncratic labour income shocks as well as permanent differences in the discount rate. Households self-insure using a risk-free government bond. Hence, our model combines a multi-region currency union model (e.g.~\cite{Gali2008-jt}; ~\cite{Nakamura2011-dw}) with a heterogeneous agent, new Keynesian model (e.g.~\cite{McKay2016-si}; ~\cite{Kaplan2018-sr}).

We calibrate the model using standard parameter values and targets as in typical new Keynesian models with heterogeneous agents. For the strength of trade linkages, we use data on shipments of goods across U.S. regions from the Commodity Flow Survey (CFS). Without


explicitly targeting the response, our model generates a local, non-durable consumption multiplier equal to 0.20 -fairly close to our estimate from the regional data. The aggregate consumption multiplier in our model is equal to 0.41 .

Our model generates positive consumption responses at the local and the aggregate levels, which is generally difficult to deliver in more standard models.\footnote{The standard neoclassical model without capital exhibits a negative aggregate consumption multiplier (and thus a less-than-one output multiplier); government spending decreases consumption due to a negative wealth effect induced by higher taxes and also due to a higher real interest rate (~\cite{Barro1984-ue}; ~\cite{Baxter1990-mo}; ~\cite{Woodford2011-jp}).} The key necessary element to generate a positive local multiplier is incomplete markets. With complete markets, any change in regional income is offset by transfers due to state-contingent claims. As a result, differences in regional consumer spending are pinned down only by differences in regional prices. And since regions with larger fiscal stimulus injections also experience higher inflation, the local consumption multiplier is negative (see, for example,~\cite{Nakamura2011-dw}; ~\cite{Farhi2012-oe}; ~\cite{Chodorow-Reich2019-zl}).

Heterogeneity is not necessary to generate a positive local or aggregate consumption response but it is important to generate substantial consumption responses consistent with the data. We show that a new Keynesian model with a representative agent in each region and incomplete markets across regions generates a positive but small local multiplier equal to 0.09 . The main reason for the larger consumption responses in our benchmark model is the substantial response of high-marginal-propensity to consume (MPC), labour-income-dependent households who experience an increase in their labour earnings. At an annual frequency, the average MPC in our model is 0.37 , which is consistent with the empirical evidence on the magnitude of non-durable consumption responses to unexpected income transfers (~\cite{Carroll2017-cp}). The combined evidence of our local consumption multiplier and existing estimates of the MPC favour a new Keynesian model with heterogeneous agents over a new Keynesian model with a representative agent.

The positive aggregate multiplier in the model also depends on our assumed passive monetary response to the fiscal stimulus. In our period of analysis, 2008-12, the Federal Reserve was at the zero lower bound (ZLB) of its policy rate and did not raise the rate in response to potential inflation pressures driven by the ARRA. Increased expected inflation in the face of an unchanged nominal rate reduces the real interest rate, leading households to increase consumption. Besides this standard effect of fiscal policy at the ZLB, a lower real interest rate also decreases the government's debt service cost which allows the government to engage in fiscal stimulus with a relatively small increase in taxes.\footnote{The decrease in the government's debt service cost is a redistribution of resources from the private to the public sector, which hurts net savers, namely wealthy, low-MPC households. In contrast, the small adjustment in taxes affects a broader group of consumers including low-income, high-MPC households. Auclert~\cite{Auclert2019-dr} analyses the redistribution channel from monetary policy to consumer spending.}

The aggregate multiplier is twice as large as the local multiplier because trade linkages propagate government spending across regions. As trade linkages become stronger, regional consumption responses comove to a larger degree and the local multiplier-which is identified by the relative cross-regional responses-decreases. We empirically validate this mechanism using trade flows from the CFS. The local multiplier is lower when estimated from regions in close trade relationships than when estimated using regions in distant trade relationships.

We show that one additional reason the aggregate multiplier is higher that the local multiplier is the passive stance of the monetary authority. If in our model the monetary policy actively responds to inflation pressures, the aggregate multiplier turns negative. Nonetheless, the local multiplier remains positive and unaffected because the monetary policy is common across regions.

Our model maps the empirical evidence on the local multiplier to an estimate for the aggregate multiplier. We show that when the model is also informed by the local multiplier, it delivers a tighter range of estimates for the aggregate multiplier relative to a model that is only informed by the MPC.

Our paper contributes to the extensive literature on the consumption multiplier. One strand of the literature estimates the aggregate consumption multiplier using aggregate vector autoregressions (VARs) (e.g.~\cite{Ramey1998-as}; Blanchard and Perotti, 2002;~\cite{Perotti2004-fu};~\cite{Barro2010-yy}). A second strand estimates multipliers using dynamic general equilibrium models (for example,~\cite{Baxter1990-mo};~\cite{Christiano2011-ro};~\cite{Drautzburg2015-ie}, to name a few).

Our paper differs from this literature on two points. First, we use cross-regional variation to identify the (local) effect of fiscal stimulus on consumer spending. With disaggregate, geographical data, we use many more observations than typically used in the VAR studies. Moreover, we can identify exogenous variation in a much broader class of government spending than time-series studies which often rely on defense spending variation. Second, we translate the cross-regional variation into an aggregate consumption response using a quantitative model. Typical dynamic equilibrium models do not rely on any cross-regional or cross-sectional evidence.

The closest paper to ours methodologically is Nakamura and Steinsson~\cite{Nakamura2011-dw}, who use cross-state evidence to analyse the output effect of defense spending. Our paper analyses the cross-regional response of consumption using detailed micro-level evidence. Additionally, Nakamura and Steinsson~\cite{Nakamura2011-dw} employ a model with complete markets and rely on nonseparable preferences between consumption and leisure to match the empirical evidence of a local output multiplier larger than one. Our model, in contrast, generates a positive local consumption multiplier due to incomplete markets.

We also contribute to the literature that uses regional variation to estimate regional effects of shocks or policies. These include work on the regional effects of house price shocks on consumer spending (~\cite{Mian2013-wr}) and the effect of unemployment insurance across regions (~\cite{HagedornUnknown-qq}). Another strand explicitly analyses the effect of fiscal stimulus on employment and income including Wilson~\cite{Wilson2012-bc}, Chodorow-Reich et al.~\cite{Chodorow-Reich2012-bo}, Conley and Dupor~\cite{Conley2013-ca}, Leduc and Wilson~\cite{Leduc2017-yn}, Serrato and Wingender~\cite{Serrato2016-wo}, for example. The above literature typically ignores general equilibrium effects. In contrast, we show how local estimates can vary from the aggregate using a general equilibrium model.

Finally, we contribute to the growing literature that incorporates household heterogeneity and incomplete markets into a new Keynesian framework. Oh and Reis~\cite{Oh2012-ax} and McKay and Reis~\cite{McKay2016-pm} study the effects of government intervention on the U.S. business cycle. Hagedorn et al.~\cite{Hagedorn2019-op}, Bhandari et al. (2018), and Auclert et al.~\cite{Auclert2018-ii} study demand shocks and fiscal policy with heterogeneity and incomplete markets. McKay et al.~\cite{McKay2016-si} and Kaplan et al.~\cite{Kaplan2018-sr} study the effects of monetary policy with heterogeneous agents. We extend this literature by incorporating multiple regions that are linked through trade, fiscal, and monetary policy. We show that, with multiple regions, the transmission of local fiscal policy depends on relative price adjustments and the strength of trade linkages, which are not considered in single-region heterogeneous agents models.

The rest of the paper is structured as follows. Section 2 describes our data. Section 3 describes our empirical specifications and documents the basic empirical patterns regarding the local response of consumer spending to local government spending. Section 4 sets up the model.

Section 5 describes our calibration and our main quantitative experiment. Section 6 analyses our results under different model specifications, and Section 7 concludes.

% Smart bibliography: Only include bibliography if standalone AND has citations
\smartbib

\end{document}
